\documentclass[11pt,letterpaper]{article}

\usepackage[margin=1in]{geometry}
\usepackage[T1]{fontenc}
\usepackage{lmodern}
\usepackage{microtype}
\usepackage{parskip}
\usepackage{enumitem}
\usepackage{booktabs}
\usepackage{longtable}
\usepackage{array}
\usepackage{titlesec}
\usepackage[hidelinks,colorlinks=true,linkcolor=black,urlcolor=black,citecolor=black]{hyperref}
\usepackage{xcolor}
\usepackage{fancyvrb}
\usepackage{fancyhdr}
\usepackage{listings}

\titleformat{\section}{\Large\bfseries}{\thesection}{1em}{}
\titleformat{\subsection}{\large\bfseries}{\thesubsection}{1em}{}
\titleformat{\subsubsection}{\normalsize\bfseries}{\thesubsubsection}{1em}{}

\pagestyle{fancy}
\fancyhf{}
\fancyhead[L]{\small Theseus --- Desktop Packaging Guide}
\fancyhead[R]{\small\thepage}
\renewcommand{\headrulewidth}{0.4pt}

\setlength{\parindent}{0pt}

\lstset{
  basicstyle=\small\ttfamily,
  breaklines=true,
  frame=single,
  backgroundcolor=\color{gray!8},
  xleftmargin=1em,
  framexleftmargin=0.5em,
  columns=fullflexible,
  keepspaces=true,
  showstringspaces=false,
  literate={~}{{\textasciitilde}}1,
}

\newcommand{\file}[1]{\texttt{#1}}
\newcommand{\cmd}[1]{\texttt{#1}}

\title{\textbf{Theseus}\\[4pt]\large Desktop Packaging Guide}
\date{April 2026}
\author{}

\begin{document}

\maketitle
\thispagestyle{empty}

\tableofcontents
\newpage

% =======================================================================
\section{Overview}
% =======================================================================

Theseus ships three desktop applications: Dialectic (PyQt6 conversation analyzer), the Noosphere CLI (epistemic engine), and the Founder Portal (Electron-wrapped Next.js). Each is distributed as a standalone installer for macOS and Windows so end users do not need Python, Node.js, or any development toolchain installed.

This guide covers how to build, sign, notarize, and package each application. It assumes you are working from a clean checkout of the Theseus monorepo on the target platform.

\subsection{Architecture at a glance}

\begin{table}[h]
\small
\begin{tabular}{p{3.2cm} p{3cm} p{4cm} p{4cm}}
\toprule
\textbf{Application} & \textbf{Packager} & \textbf{macOS output} & \textbf{Windows output} \\
\midrule
Dialectic           & PyInstaller    & \file{.app} in DMG          & \file{.exe} via NSIS installer \\
Noosphere CLI       & PyInstaller    & \file{tar.gz} archive       & \file{.exe} via NSIS installer \\
Founder Portal      & electron-builder & \file{.app} in DMG        & \file{.exe} via NSIS installer \\
\bottomrule
\end{tabular}
\end{table}

% =======================================================================
\section{Prerequisites}
% =======================================================================

\subsection{All platforms}

\begin{itemize}[leftmargin=1.5em,itemsep=0.2em]
  \item Git checkout of the Theseus monorepo.
  \item Python 3.11+ with pip (for Dialectic and Noosphere).
  \item Node.js 20+ with npm (for Founder Portal).
  \item Pillow (Python) for icon generation: \cmd{pip install pillow}.
\end{itemize}

\subsection{macOS-specific}

\begin{itemize}[leftmargin=1.5em,itemsep=0.2em]
  \item Xcode Command Line Tools (\cmd{xcode-select --install}).
  \item \cmd{create-dmg} (optional, installed via \cmd{brew install create-dmg}; the build script falls back to \cmd{hdiutil} if absent).
  \item For code signing: an Apple Developer ID Application certificate in your Keychain, identified by its Common Name (e.g., \cmd{Developer ID Application: Theseus Inc (TEAMID)}).
  \item For notarization: an Apple ID enrolled in the Apple Developer Program, with an app-specific password generated at \cmd{appleid.apple.com}.
\end{itemize}

\subsection{Windows-specific}

\begin{itemize}[leftmargin=1.5em,itemsep=0.2em]
  \item Windows SDK (provides \cmd{signtool.exe} for Authenticode signing).
  \item NSIS 3.x (Nullsoft Scriptable Install System) --- download from \cmd{nsis.sourceforge.io} and ensure \cmd{makensis} is on PATH.
  \item For code signing: a \file{.pfx} code signing certificate and its password.
\end{itemize}

% =======================================================================
\section{Dialectic (PyQt6 desktop app)}
% =======================================================================

\subsection{Key files}

\begin{table}[h]
\small
\begin{tabular}{p{5.5cm} p{9cm}}
\toprule
\textbf{File} & \textbf{Purpose} \\
\midrule
\file{dialectic/dialectic.spec}             & PyInstaller spec --- defines entry point, hidden imports, data bundles, exclusions, platform targets. \\
\file{dialectic/dialectic/resources.py}     & Frozen-mode path resolver. Returns correct asset/data paths whether running from source or PyInstaller bundle. \\
\file{dialectic/dialectic/updater.py}       & Background update checker. Polls a JSON manifest URL; notifies user in the Qt dashboard. \\
\file{dialectic/scripts/generate\_icons.py} & Converts \file{assets/icon.png} to \file{.icns} (macOS) and \file{.ico} (Windows). \\
\file{dialectic/scripts/build\_macos.sh}    & Full macOS build: icons, PyInstaller, DMG creation, optional signing and notarization. \\
\file{dialectic/scripts/build\_windows.ps1} & Full Windows build: icons, PyInstaller. NSIS installer is a separate step. \\
\file{dialectic/installer/dialectic.nsi}    & NSIS installer script: install directory, Start Menu shortcuts, Desktop shortcut, uninstaller, registry. \\
\file{dialectic/installer/dmg-settings.json}& DMG layout: drag-to-Applications, icon positions, background color. \\
\bottomrule
\end{tabular}
\end{table}

\subsection{Building on macOS}

\begin{lstlisting}
cd dialectic

# 1. Install Python dependencies
pip install -r requirements.txt
pip install pyinstaller pillow

# 2. Run the full build (icons + PyInstaller + DMG)
bash scripts/build_macos.sh
\end{lstlisting}

The script performs these steps internally:
\begin{enumerate}[leftmargin=1.5em,itemsep=0.1em]
  \item Runs \file{generate\_icons.py} to produce \file{assets/Dialectic.icns}.
  \item Invokes \cmd{pyinstaller dialectic.spec --noconfirm --clean}.
  \item Reads the version from \file{pyproject.toml}.
  \item Creates \file{dist/Dialectic-\{version\}.dmg} using \cmd{create-dmg} (with drag-to-Applications layout) or \cmd{hdiutil} as fallback.
\end{enumerate}

Output: \file{dialectic/dist/Dialectic.app} and \file{dialectic/dist/Dialectic-0.1.0.dmg}.

\textbf{Signing and notarization.} The build script reads environment variables:
\begin{lstlisting}
export APPLE_SIGNING_IDENTITY="Developer ID Application: ..."
export APPLE_ID="your@apple.id"
export APPLE_TEAM_ID="TEAMID"
export APPLE_APP_PASSWORD="xxxx-xxxx-xxxx-xxxx"

bash scripts/build_macos.sh
\end{lstlisting}

If these are set, the script calls \file{scripts/codesign\_macos.sh} (deep-signs all binaries with hardened runtime) and \file{scripts/notarize\_macos.sh} (submits the DMG to Apple, waits up to 600s, staples the ticket). If the variables are absent, both steps are skipped with a warning and the build still succeeds unsigned.

\subsection{Building on Windows}

\begin{lstlisting}
cd dialectic

# 1. Install dependencies
pip install -r requirements.txt
pip install pyinstaller pillow

# 2. Build with PyInstaller
powershell -File scripts\build_windows.ps1

# 3. Create the NSIS installer
makensis installer\dialectic.nsi
\end{lstlisting}

Output: \file{dialectic\textbackslash dist\textbackslash Dialectic\textbackslash} (directory with \file{Dialectic.exe}) and \file{dialectic\textbackslash dist\textbackslash Dialectic-Setup.exe} (installer).

The NSIS installer installs to \file{C:\textbackslash Program Files\textbackslash Theseus\textbackslash Dialectic}, creates Start Menu and Desktop shortcuts, registers an uninstaller in Windows Add/Remove Programs, and writes registry entries for display name, publisher, and icon.

\textbf{Signing.} After building but before distribution:
\begin{lstlisting}
powershell -File ..\scripts\codesign_windows.ps1 `
  -Target dist\Dialectic `
  -CertPath C:\path\to\cert.pfx `
  -CertPassword "your-password"
\end{lstlisting}

This signs every \file{.exe} and \file{.dll} in the output directory with SHA-256 and a DigiCert timestamp.

\subsection{Bundle size and ML models}

The Dialectic bundle is large (${\sim}$2--3\,GB) because it includes PyTorch, faster-whisper (CTranslate2), sentence-transformers, and their transitive native libraries. The PyInstaller spec uses single-directory mode (not single-file) because extracting a multi-gigabyte archive on every launch would be impractical.

The spec excludes packages not needed at runtime: \cmd{tensorflow}, \cmd{tensorboard}, \cmd{jupyter}, \cmd{IPython}, \cmd{pytest}, and \cmd{matplotlib}. If bundle size becomes critical, the two largest contributors are PyTorch (${\sim}$800\,MB) and the CTranslate2 Whisper model (${\sim}$500\,MB). A future optimization could download models on first launch rather than bundling them.

\subsection{Update mechanism}

\file{dialectic/dialectic/updater.py} runs a background thread on application launch that fetches a JSON manifest from a configurable URL (default: \cmd{https://releases.theseus.co/dialectic/latest.json}). The manifest format is:
\begin{lstlisting}
{
  "version": "0.2.0",
  "download_url": "https://...",
  "release_notes": "Bug fixes and ..."
}
\end{lstlisting}

If the manifest version is newer than the running version, the callback notifies the user in the Qt dashboard. The check is non-blocking (10s timeout), daemon-threaded, and silently no-ops on network errors. There is no silent auto-install; the user must download and install manually. To host the manifest, place the JSON file at the configured URL and update it on each release.

% =======================================================================
\section{Noosphere CLI}
% =======================================================================

\subsection{Key files}

\begin{table}[h]
\small
\begin{tabular}{p{5.8cm} p{8.7cm}}
\toprule
\textbf{File} & \textbf{Purpose} \\
\midrule
\file{noosphere/noosphere.spec}                 & PyInstaller spec for the CLI binary. \\
\file{noosphere/noosphere/frozen\_support.py}   & Frozen-mode path resolver for bundle dir, data dir, Alembic dir. \\
\file{noosphere/noosphere/updater.py}           & Background update checker with Rich terminal output. \\
\file{noosphere/scripts/build\_macos.sh}        & macOS build: PyInstaller + tar.gz archive. \\
\file{noosphere/scripts/build\_windows.ps1}     & Windows build: PyInstaller. \\
\file{noosphere/installer/noosphere.nsi}        & NSIS installer: installs to Program Files and adds to PATH. \\
\file{noosphere/scripts/build\_hooks/hook-noosphere.py} & PyInstaller runtime hook setting \cmd{NOOSPHERE\_FROZEN=1}. \\
\bottomrule
\end{tabular}
\end{table}

\subsection{Building on macOS}

\begin{lstlisting}
cd noosphere
pip install -r requirements.txt
pip install pyinstaller
bash scripts/build_macos.sh
\end{lstlisting}

Output: \file{noosphere/dist/noosphere-0.1.0-macos.tar.gz}. Users extract this archive and add the \file{noosphere/} directory to their PATH.

Unlike Dialectic, the Noosphere CLI does not produce a \file{.app} bundle because it is a terminal tool, not a GUI application. A DMG option exists (\file{installer/dmg-settings.json} with a drag-to-\file{/usr/local/bin} layout) but the tar.gz is the primary distribution.

Signing follows the same pattern as Dialectic: set \cmd{APPLE\_SIGNING\_IDENTITY} and the notarization variables, and the build script calls the shared signing scripts.

\subsection{Building on Windows}

\begin{lstlisting}
cd noosphere
pip install -r requirements.txt
pip install pyinstaller
powershell -File scripts\build_windows.ps1
makensis installer\noosphere.nsi
\end{lstlisting}

The NSIS installer adds the install directory to the system PATH so that \cmd{noosphere} is available from any terminal.

\subsection{Alembic migrations in frozen mode}

The spec bundles \file{alembic/} and \file{alembic.ini} as data files. The \file{frozen\_support.py} module provides \cmd{alembic\_dir()} and \cmd{alembic\_ini()} which resolve to the bundled copies when frozen. Any code that runs migrations must use these paths rather than assuming the working directory contains Alembic files.

\subsection{Data directory}

The frozen CLI writes user data (SQLite database, graph, embeddings, synthesis output) to a platform-specific directory:

\begin{itemize}[leftmargin=1.5em,itemsep=0.1em]
  \item macOS: \file{\textasciitilde/Library/Application Support/Noosphere}
  \item Windows: \file{\%APPDATA\%\textbackslash Noosphere}
  \item Linux: \file{\textasciitilde/.noosphere}
\end{itemize}

The \cmd{THESEUS\_DATA\_DIR} environment variable overrides this default. The directory is created automatically on first run.

ML model weights (sentence-transformers, spaCy) are \textbf{not} bundled --- they are too large (${\sim}$2\,GB). On first run of any command that requires embeddings, the models download to the standard Hugging Face cache directory. Users behind firewalls must pre-download models or configure \cmd{HF\_HOME} to point to a local mirror.

% =======================================================================
\section{Founder Portal (Electron)}
% =======================================================================

\subsection{Key files}

\begin{table}[h]
\small
\begin{tabular}{p{6cm} p{8.5cm}}
\toprule
\textbf{File} & \textbf{Purpose} \\
\midrule
\file{founder-portal/desktop/main.ts}          & Electron main process: lifecycle, window, migrations. \\
\file{founder-portal/desktop/next-server.ts}    & Starts the Next.js standalone server as a child process. \\
\file{founder-portal/desktop/db-path.ts}        & Resolves SQLite database to platform app-data directory. \\
\file{founder-portal/desktop/preload.ts}        & Context bridge (platform info, deep-link IPC). \\
\file{founder-portal/desktop/updater.ts}        & \cmd{electron-updater} integration for auto-update. \\
\file{founder-portal/desktop/dev.ts}            & Dev mode: starts Next.js dev server + Electron. \\
\file{founder-portal/electron-builder.yml}      & Build config: DMG, NSIS, signing, notarization, publish. \\
\file{founder-portal/assets/entitlements.mac.plist} & macOS entitlements for hardened runtime. \\
\file{founder-portal/scripts/generate\_icons.js}    & Generates \file{.icns} and \file{.ico} from source PNG. \\
\file{founder-portal/scripts/build\_desktop\_macos.sh} & Full macOS desktop build. \\
\file{founder-portal/scripts/build\_desktop\_windows.ps1} & Full Windows desktop build. \\
\bottomrule
\end{tabular}
\end{table}

\subsection{How the Electron wrapper works}

The Electron main process performs this startup sequence:

\begin{enumerate}[leftmargin=1.5em,itemsep=0.1em]
  \item Acquires a single-instance lock (prevents duplicate launches).
  \item Resolves the SQLite database path via \file{db-path.ts}: \file{userData/data/founder-portal.db} (platform-specific --- \file{\textasciitilde/Library/Application Support/Theseus Founder Portal/} on macOS, \file{\%APPDATA\%/Theseus Founder Portal/} on Windows).
  \item If the database file does not exist (first launch), runs Prisma migrations or copies a pre-migrated seed database.
  \item Calls \file{next-server.ts} which \cmd{fork}s the Next.js standalone \file{server.js} on a free port (starting from 3000). The function polls the port with exponential backoff (max 30s timeout) to confirm the server is ready.
  \item Creates a \cmd{BrowserWindow} (1400$\times$900, min 1024$\times$700, context isolation enabled, node integration disabled, \cmd{titleBarStyle: hiddenInset} on macOS).
  \item Loads \cmd{http://127.0.0.1:\{port\}}.
  \item Intercepts external link navigation and opens in the system browser.
  \item Initializes the auto-updater (wrapped in try/catch so missing GitHub Release config does not crash the app).
\end{enumerate}

On quit, the main process sends SIGTERM to the Next.js child process, waits 5s, then SIGKILL if it has not exited.

\subsection{Building on macOS}

\begin{lstlisting}
cd founder-portal
npm ci
bash scripts/build_desktop_macos.sh
\end{lstlisting}

The script runs: icon generation, \cmd{npm run build} (Next.js), TypeScript compilation of the \file{desktop/} modules, and \cmd{electron-builder --mac --config electron-builder.yml}.

Output: \file{founder-portal/dist-electron/Theseus Founder Portal-\{version\}.dmg}.

\textbf{Code signing.} electron-builder handles signing automatically when these environment variables are set:
\begin{lstlisting}
export CSC_LINK="/path/to/cert.p12"     # or base64-encoded
export CSC_KEY_PASSWORD="cert-password"
\end{lstlisting}

The \file{electron-builder.yml} specifies \cmd{hardenedRuntime: true} and references \file{assets/entitlements.mac.plist} which grants JIT execution, unsigned executable memory (required by V8), dynamic linker environment variables, network client access, and user-selected file read/write.

\textbf{Notarization.} electron-builder submits to Apple automatically when the following are also set:
\begin{lstlisting}
export APPLE_ID="your@apple.id"
export APPLE_ID_PASSWORD="app-specific-password"
export APPLE_TEAM_ID="TEAMID"
\end{lstlisting}

\subsection{Building on Windows}

\begin{lstlisting}
cd founder-portal
npm ci
powershell -File scripts\build_desktop_windows.ps1
\end{lstlisting}

Output: \file{founder-portal\textbackslash dist-electron\textbackslash Theseus Founder Portal Setup \{version\}.exe}.

\textbf{Code signing.} Set environment variables before building:
\begin{lstlisting}
$env:WIN_CSC_LINK = "C:\path\to\cert.pfx"
$env:WIN_CSC_KEY_PASSWORD = "cert-password"
\end{lstlisting}

electron-builder signs the executable and the NSIS installer wrapper automatically.

\subsection{Auto-update}

The Founder Portal uses \cmd{electron-updater} for automatic updates via GitHub Releases. The \file{electron-builder.yml} contains a \cmd{publish} block:
\begin{lstlisting}
publish:
  provider: github
  owner: OWNER_PLACEHOLDER
  repo: REPO_PLACEHOLDER
\end{lstlisting}

Replace the placeholders with the actual GitHub repository coordinates. When a new release is published (by the CI/CD release workflow or manually), electron-updater detects it on the next launch, downloads the update in the background, and prompts the user to restart. The update flow is: check $\to$ download $\to$ dialog (``Restart'' / ``Later'') $\to$ quit-and-install on confirm.

The updater is initialized in \file{desktop/main.ts} after the window loads, wrapped in try/catch. If the publish config is not set or the GitHub token is missing, the updater silently fails without affecting app functionality.

\subsection{Development workflow}

For local iteration without building a full package:
\begin{lstlisting}
cd founder-portal
npm ci
npm run desktop:dev
\end{lstlisting}

This runs \file{desktop/dev.ts}, which starts the Next.js dev server on port 3000 (with hot reload), waits for it to be ready, then launches Electron pointing at \cmd{localhost:3000}. Changes to Next.js pages and components are reflected immediately; changes to \file{desktop/*.ts} require restarting the command.

% =======================================================================
\section{Shared signing infrastructure}
% =======================================================================

Three scripts in the repository root \file{scripts/} directory provide signing capabilities shared by all applications. Each is designed for both local use and CI integration: they accept credentials via arguments or environment variables, and exit 0 with a warning when credentials are not provided.

\subsection{macOS code signing}

\file{scripts/codesign\_macos.sh} deep-signs a \file{.app} bundle or PyInstaller output directory:

\begin{enumerate}[leftmargin=1.5em,itemsep=0.1em]
  \item Finds all nested binaries (\file{.dylib}, \file{.so}, executables) via \cmd{find}.
  \item Signs each with \cmd{codesign --force --options runtime --timestamp --sign \$IDENTITY}.
  \item Signs the top-level target with \cmd{--deep}.
  \item Verifies with \cmd{codesign --verify --deep --strict}.
\end{enumerate}

Usage: \cmd{bash scripts/codesign\_macos.sh <path> <identity>}. Pass \cmd{skip} as the identity to skip signing.

\subsection{macOS notarization}

\file{scripts/notarize\_macos.sh} submits a DMG, ZIP, or app bundle to Apple:

\begin{enumerate}[leftmargin=1.5em,itemsep=0.1em]
  \item Calls \cmd{xcrun notarytool submit} with Apple ID, Team ID, and app-specific password.
  \item Waits up to 600 seconds for Apple to process.
  \item Staples the ticket to the artifact via \cmd{xcrun stapler staple}.
\end{enumerate}

Usage: \cmd{bash scripts/notarize\_macos.sh <artifact> <apple-id> <team-id> <password>}. Skips gracefully when any credential is empty.

\subsection{Windows Authenticode signing}

\file{scripts/codesign\_windows.ps1} signs all executables and DLLs in a target directory:

\begin{enumerate}[leftmargin=1.5em,itemsep=0.1em]
  \item Locates \cmd{signtool.exe} from the Windows SDK installation.
  \item Iterates over all \file{.exe} and \file{.dll} files recursively.
  \item Signs each with SHA-256 digest and DigiCert timestamp server.
\end{enumerate}

Usage: \cmd{.\textbackslash scripts\textbackslash codesign\_windows.ps1 -Target <dir> -CertPath <pfx> -CertPassword <pw>}. Skips gracefully when the certificate path is empty or the file does not exist.

% =======================================================================
\section{Troubleshooting}
% =======================================================================

\textbf{PyInstaller ``hidden import'' errors.} If the built app crashes with \cmd{ModuleNotFoundError}, add the missing module to the \cmd{hiddenimports} list in the \file{.spec} file. Common culprits: \cmd{sklearn.utils.\_cython\_blas}, \cmd{scipy.special.\_cdflib}, \cmd{torch.\_C}. Run \cmd{pyinstaller --debug=imports} for a full import trace.

\textbf{macOS notarization rejected.} Apple rejects binaries that are not signed with hardened runtime or that contain unsigned nested binaries. Ensure the \file{codesign\_macos.sh} script runs \emph{before} creating the DMG, and that every \file{.dylib} inside the bundle is individually signed. Check rejection details with \cmd{xcrun notarytool log <submission-id>}.

\textbf{Electron builder fails on native modules.} \cmd{better-sqlite3} is a native Node module that must be rebuilt for Electron's Node version. electron-builder handles this via \cmd{electronRebuild}, but if it fails, run \cmd{npx electron-rebuild -f -w better-sqlite3} manually before building.

\textbf{Windows NSIS ``file not found.''} The \file{.nsi} script references \file{..\textbackslash dist\textbackslash Dialectic\textbackslash *.*} --- ensure PyInstaller has been run first and the output is in \file{dist\textbackslash Dialectic\textbackslash}.

\textbf{Bundle size too large.} The two largest contributors in Python bundles are PyTorch and CTranslate2. To reduce size: (a) use \cmd{torch} CPU-only builds (\cmd{pip install torch --index-url https://download.pytorch.org/whl/cpu}), (b) strip debug symbols (\cmd{strip -x} on macOS, \cmd{strip --strip-unneeded} on Linux), (c) exclude unused torch backends via spec \cmd{excludes}.

\end{document}
